% ==============================================================
% contents.tex — write your paper here
% ==============================================================
% AUTHORS. This is the file you edit day to day, together with
% article.toml (title, abstracts, keywords, highlights),
% authors.toml and references.bib.
%
% It begins at the first section: the whole cover -- journal
% masthead, title, author list, abstracts and the sidebar -- is
% generated from the TOML files, so there is no front matter to
% maintain here.
%
% It ships as a working skeleton of a real article: Introduction,
% Methods, Results, Conclusions and two appendices. The Methods
% section doubles as the manual -- its subsections demonstrate
% every feature the template provides (citations, equations,
% figures, tables). Replace the blind text with your own writing
% and delete whatever you don't need.
% ==============================================================

% ==============================================================
% 1. INTRODUCTION
% ==============================================================
\section{Introduction} \label{sec:introduction}

Between the end of April and the beginning of May 2024, heavy
rains caused extraordinary flooding in southern Brazil. The
largest floods and landslides ever observed occurred along
several rivers in the state of Rio Grande do Sul, such as the
Jacu\'i, Taquari, Ca\'i, Sinos and Gua\'iba
\parencite{nytimes2024}. The Patos Lagoon, which ultimately
receives the water of all those rivers, also reached the highest
levels ever recorded. In the state's capital, Porto Alegre, the
Gua\'iba flood exceeded the maximum level observed during the
traumatic flood of 1941 \parencite{possa2022b,silveira2023}, with
even more devastating consequences because the flood protection
system, built in the 1970 decade, was unable to prevent the water
from entering the city.

The hydrological disaster occurred in four forms: (1) mass
movements in the mountains, (2) flash floods in rivers in the
transition between the mountains and the plains, (3) relatively
slower and long-lasting flooding along the densely populated
areas close to the lowland rivers, where the failure of the flood
protection systems intensified the impacts, and (4) slow flooding
throughout the margins of the Patos Lagoon.

The exceptional hydrological disaster of April--May 2024 in
southern Brazil not only marks new records of flooding in several
rivers in the region; it is also the most extensive landslide
event ever observed in Brazil
\parencite{cemaden2024a,andrades2024}. Hundreds of thousands of
people were directly affected by the floods, which inundated more
than 15,000~km\textsuperscript{2}. Fatal victims are at least
183, and 27 people are still reported missing, five months after
the disaster \parencite{rs2024d}.

The 2024 extreme hydrological event occurred a few months after
other noteworthy rainfall-related disasters in Brazil
\parencite{marengo2023a,marengo2023b,marengo2024}, and only eight
months after a catastrophic flood hit the same southern Brazilian
region in September 2023 \parencite{alvala2024,moraes2024},
suggesting that the region is passing through a particular
flood-rich period. This recent flood-rich period in southern
Brazil seems to be in accordance with already reported trends in
observed rainfall and floods during the last decades
\parencite{chagas2022} and with projections of possible
hydrological impacts of climate change in southern Brazil for the
twenty-first century \parencite{breda2023,hodnebrog2022}.

Given its profound, widespread and long-lasting impacts, the
disaster of April--May 2024 can be considered the biggest natural
hydrological disaster in the history of Brazil. It also gives an
example of an event that can be called an ``impossible flood'',
as suggested by \textcite{montanari2024}, because it exceeded the
expectations based on historical experience, and because it
exposed the consequences of the occurrence of an event whose
probability of occurrence was considered too small to act, even
by the flood managers \parencite{allasia2015}.

The objective of this study is to document this exceptional
hydrological event and the disaster that followed. We aimed to
describe the process chain of the April--May 2024 hydrological
disaster in South Brazil, from causes to impacts, focusing on
what made the impacts so exceptional.

In our understanding, adequately documenting extraordinary
hydrological events helps society learn from those events, by
understanding patterns and causes. This learning process helps
reduce vulnerability and increase resilience to future disasters
\parencite{zanandrea2023}. Other examples of papers that describe
extraordinary floods are given by \textcite{gaume2004},
\textcite{dimitriou2024}, \textcite{mohr2023},
\textcite{nanditha2023} and \textcite{kreibich2017}. In the
Brazilian context, accounts of other extraordinary hydrological
events have been recently organized by \textcite{zanandrea2023}.
We also understand that a detailed description and study of such
an extraordinary flood event is not of exclusive interest to the
affected region, because the lessons that can be drawn from it
are not limited to the basin or river in which it occurred but
can be extremely relevant to an entire large region, as suggested
by \textcite{merz2024}.

The paper is structured as follows: \Cref{sec:methods} describes
the data sets and methods used. \Cref{sec:results} describes the
event in terms of meteorological and hydrological aspects and its
impacts. Finally, \Cref{sec:conclusions} discusses and summarizes
the main results and draws conclusions.

% ==============================================================
% 2. METHODS
% ==============================================================
% This section doubles as the template's manual: each subsection
% demonstrates one feature. Keep the structure, replace the
% content.
\section{Methods} \label{sec:methods}

\lipsum[4]

\subsection{Citing sources} \label{sec:citations}

References live in \texttt{references.bib} and are formatted in
APA style with author--year citations. There are two forms, and
the choice is grammatical rather than cosmetic.

Use \verb|\parencite{key}| when the citation supports a statement
but is not part of the sentence, so it belongs in parentheses:
the Patos Lagoon is the largest choked lagoon in South America
\parencite{chagas2022}. Use \verb|\textcite{key}| when the
authors \emph{are} the subject of the sentence, so their names
need to read as words: \textcite{montanari2024} calls such an
event an ``impossible flood''.

Both take several keys at once, separated by commas ---
\verb|\parencite{a,b}| --- and biblatex orders and punctuates the
group for you \parencite{marengo2023a,marengo2023b,marengo2024}.
Every citation is a clickable cross-reference to its entry in the
reference list.

Cross-references to document elements use \verb|\Cref{label}|,
which supplies the word as well as the number
(\Cref{sec:introduction}), and \verb|\cref{label}| inside a
sentence, as in \cref{sec:tables} below. Give every figure,
table, equation and section a label and never type a number by
hand.

\lipsum[5-8]

\subsection{Mathematical features} \label{sec:equations}

Inline mathematics is written between dollar signs, as in the
continuity relation $Q = v\,A$, where $Q$ is discharge, $v$ mean
velocity and $A$ the cross-sectional area. Keep short expressions
inline; give anything you need to refer to a display of its own.

Numbered display equations go in an \texttt{equation} environment
with a label, so they can be cross-referenced:

\begin{equation}
    \label{eq:manning}
    v = \frac{1}{n}\,R^{2/3}\,S^{1/2}.
\end{equation}

\Cref{eq:manning} is referenced by label, so its number stays
correct however many equations are inserted before it. Use
\verb|\Cref| at the start of a sentence and \verb|\cref| inside
one.

An unnumbered display uses the starred form, for a step that
needs room but never gets referred to:

\begin{equation*}
    R = \frac{A}{P}.
\end{equation*}

Related equations that belong together can share a number and be
lettered individually with \texttt{subequations}:

\begin{subequations}
    \label{eq:balance}
    \begin{align}
        \frac{\partial A}{\partial t} + \frac{\partial Q}{\partial x} &= q_{\ell}, \label{eq:balance:mass} \\
        \frac{\partial Q}{\partial t} + \frac{\partial}{\partial x}\!\left(\frac{Q^{2}}{A}\right) &= gA\left(S_{0} - S_{f}\right). \label{eq:balance:momentum}
    \end{align}
\end{subequations}

\Cref{eq:balance} is the pair; \cref{eq:balance:mass} and
\cref{eq:balance:momentum} are its parts, numbered \textbf{a} and
\textbf{b} automatically.

Systems are aligned on a chosen character with \texttt{align},
one \verb|&| per alignment point:

\begin{align}
    \label{eq:system}
    S_{f}   &= \frac{n^{2} v^{2}}{R^{4/3}}, &
    F_{r}   &= \frac{v}{\sqrt{g\,h}}, \\
    \tau_0  &= \rho\,g\,R\,S_{f},           &
    \mathrm{Re} &= \frac{v\,R}{\nu}. \nonumber
\end{align}

Symbols take the usual variants: Greek letters $\alpha$, $\beta$,
$\lambda$, $\mu$, $\sigma$, $\Delta$, $\Omega$; decorated forms
$\bar{x}$, $\hat{y}$, $\tilde{z}$, $\dot{Q}$, $\vec{u}$; calligraphic
$\mathcal{R}$, blackboard $\mathbb{R}$, bold $\mathbf{A}$ and
roman subscripts such as $Q_{\mathrm{peak}}$. Operators and
relations behave as expected: $\sum$, $\prod$, $\int$,
$\partial$, $\nabla$, $\leq$, $\approx$, $\propto$, $\in$.

Theorem-like environments are available for formal statements,
and number independently of the equations:

\begin{definition}[Specific energy]
    \label{def:energy}
    For a channel of depth $h$ carrying mean velocity $v$, the
    specific energy is $E = h + v^{2}/2g$.
\end{definition}

\begin{theorem}[Critical flow]
    \label{thm:critical}
    Specific energy is minimised when $F_{r} = 1$.
\end{theorem}

\begin{proof}
    Differentiate $E$ from \cref{def:energy} with respect to $h$
    at fixed discharge and set $\mathrm{d}E/\mathrm{d}h = 0$;
    the result rearranges to $v^{2} = g\,h$, which is
    $F_{r} = 1$ by \cref{eq:system}.
\end{proof}

\begin{remark}
    \texttt{lemma}, \texttt{corollary}, \texttt{example} and
    \texttt{remark} are defined alongside these, and all follow
    the document language.
\end{remark}

\lipsum[5-8]

\subsection{Developing figures} \label{sec:figures}

Figures come in three sizes, each a single macro. Give the image
file's bare name --- no folder, no extension --- and the template
looks it up in \texttt{figures/}, preferring \texttt{.pdf}, then
\texttt{.png}, then \texttt{.jpeg}, then \texttt{.jpg}:

\begin{itemize}
  \item \verb|\FigureM{name}{caption}{label}| --- text width,
        as in \Cref{fig:forecast}.
  \item \verb|\FigureS{name}{caption}{label}| --- small, with the
        caption set beside the image and spilling into the
        sidebar column, as in \Cref{fig:stage}. Pass a second
        optional argument to change its width, for example
        \verb|\FigureS[t][0.5]{...}|; the default is 0.67 of the
        text width.
  \item \verb|\FigureL{name}{caption}{label}| --- full width,
        margin to margin, as in \Cref{fig:damage}.
\end{itemize}

Every one of them takes an optional float position first, so
\verb|\FigureM[b]{...}| sends that figure to the bottom of a
page. The default is \texttt{t}, the top.

\FigureM{figure_05}{%
  Seven-day accumulated precipitation forecasts from the ECMWF
  and GFS global models for the river basins of Rio Grande do
  Sul. Darker shading indicates higher forecast accumulation;
  both models anticipated the heavy rainfall that preceded the
  flood.
  \;\textbf{a}\;---\;ECMWF, forecast issued 27 April.
  \;\textbf{b}\;---\;GFS, forecast issued 27 April.
  \;\textbf{c}\;---\;ECMWF, forecast issued 29 April.
  \;\textbf{d}\;---\;GFS, forecast issued 29 April.%
}{fig:forecast}

\lipsum[5-8]

\FigureS{figure_15}{%
  River stage (water level) from 1 May to 3 June 2024 at six
  gauging stations across the Rio Grande do Sul lowlands and
  coastal lagoon system (Arambar\'e, S\~ao Louren\c{c}o,
  Laranjal, Rio Grande, Pelotas Porto and CATSUL), showing the
  sustained high-water period that followed the flood event.%
}{fig:stage}

\lipsum[5-8]

\FigureL{figure_07}{%
  Aerial and ground documentation of slope failures and
  structural damage following the April--May 2024 event.
  \;\textbf{a}\;---\;debris-flow scars draining into a stream channel.
  \;\textbf{b}\;---\;hillslope cracking upstream of a residential area.
  \;\textbf{c}\;---\;a large landslide scar above a rural property.
  \;\textbf{d}\;---\;structural cracking in a house wall.
  \;\textbf{e}\;---\;a building destroyed by a landslide.%
}{fig:damage}

\lipsum[5-8]

If an image file cannot be found, the build still succeeds and a
placeholder appears in its place, naming the file it looked for
--- see \Cref{fig:missing}. A missing or misspelled image is
therefore visible in the PDF instead of stopping the compile.
Delete this example from your own paper.

\FigureM{figure_99}{%
  Placeholder shown when an image file cannot be found in
  \texttt{figures/}. The box keeps the golden ratio of the
  requested width.%
}{fig:missing}

\lipsum[5-8]

\subsection{Presenting tables} \label{sec:tables}

Tables come in two sizes and are written as environments, so the
tabular material goes in the body and the caption is supplied as
an argument. The caption is placed above the table, following the
usual convention:

\begin{itemize}
  \item \verb|\begin{TableM}{caption}{label}| --- text width,
        as in \Cref{tbl:stations}.
  \item \verb|\begin{TableL}{caption}{label}| --- full width,
        margin to margin, for tables with many columns, as in
        \Cref{tbl:events}.
\end{itemize}

Both accept the same optional position argument as the figure
macros, defaulting to \texttt{t}. Use \texttt{booktabs} rules
(\verb|\toprule|, \verb|\midrule|, \verb|\bottomrule|) and avoid
vertical lines. Footnotes to a table go in a \verb|tablenotes|
block at the end of the environment, keyed with \verb|\tnote{a}|
in the cell they annotate.

\begin{TableM}{%
  Gauging stations used in this study, with drainage area and the
  peak stage recorded during the event.%
}{tbl:stations}
\scriptsize
\sffamily
% tabularx with L columns fills \linewidth exactly, which keeps
% the caption and notes flush with the text column.
\begin{tabularx}{\linewidth}{LLLL}
\toprule
\textbf{Station} & \textbf{River} & \textbf{Area (km\textsuperscript{2})} & \textbf{Peak (m)} \\
\midrule
Arambaré          & Patos Lagoon  & 180{,}000 & 2.6\tnote{a} \\
São Lourenço      & Patos Lagoon  & 180{,}000 & 2.8 \\
Laranjal          & Patos Lagoon  & 180{,}000 & 2.8 \\
Pelotas Porto     & São Gonçalo   &  4{,}800  & 3.0 \\
\bottomrule
\end{tabularx}
\begin{tablenotes}
  \item[\textit{a}] Notes go in a \texttt{tablenotes} block and
  wrap to the width of the table above them.
\end{tablenotes}
\end{TableM}

\lipsum[11-13]

The large variant spans the full width, which makes room for
tables that would otherwise be cramped in the text column.



\begin{TableL}{%
  Summary of the sub-basins affected by the April--May 2024
  event, with rainfall totals, peak discharge and the recurrence
  interval estimated for each.%
}{tbl:events}
\scriptsize
\sffamily
\begin{tabularx}{\linewidth}{LLLLLL}
\toprule
\textbf{Sub-basin} & \textbf{Area (km\textsuperscript{2})} &
\textbf{Rainfall (mm)} & \textbf{Peak (m\textsuperscript{3}/s)} &
\textbf{Date} & \textbf{Return period (years)} \\
\midrule
Taquari     &  26{,}500 & 520 & 14{,}000 & 2024-05-02 & $>$100 \\
Ca\'i       &   4{,}900 & 480 &  5{,}200 & 2024-05-01 & $>$100 \\
Sinos       &   3{,}700 & 460 &  4{,}100 & 2024-05-03 & $>$100 \\
Jacu\'i     &  71{,}600 & 410 & 19{,}500 & 2024-05-04 & $>$100\tnote{a} \\
Gravata\'i  &   2{,}000 & 390 &    980   & 2024-05-05 &    50  \\
Taquari     &  26{,}500 & 520 & 14{,}000 & 2024-05-02 & $>$100 \\
Ca\'i       &   4{,}900 & 480 &  5{,}200 & 2024-05-01 & $>$100 \\
Sinos       &   3{,}700 & 460 &  4{,}100 & 2024-05-03 & $>$100 \\
Jacu\'i     &  71{,}600 & 410 & 19{,}500 & 2024-05-04 & $>$100\tnote{a} \\
Gravata\'i  &   2{,}000 & 390 &    980   & 2024-05-05 &    50  \\
Taquari     &  26{,}500 & 520 & 14{,}000 & 2024-05-02 & $>$100 \\
Ca\'i       &   4{,}900 & 480 &  5{,}200 & 2024-05-01 & $>$100 \\
Sinos       &   3{,}700 & 460 &  4{,}100 & 2024-05-03 & $>$100 \\
Jacu\'i     &  71{,}600 & 410 & 19{,}500 & 2024-05-04 & $>$100\tnote{a} \\
Gravata\'i  &   2{,}000 & 390 &    980   & 2024-05-05 &    50  \\
\bottomrule
\end{tabularx}
\begin{tablenotes}
  \item[\textit{a}] Return periods above 100 years are reported as
  a lower bound, the record being too short to resolve them.
\end{tablenotes}
\end{TableL}



\lipsum[14]

% ==============================================================
% 3. RESULTS
% ==============================================================
\section{Results} \label{sec:results}

\lipsum[15-18]

% ==============================================================
% 4. CONCLUSIONS
% ==============================================================
\section{Conclusions} \label{sec:conclusions}

\lipsum[19-23]

